% ------------------------------------------------------------------------------
% 9.2. DISEÑO DE PRUEBAS
% ------------------------------------------------------------------------------
\subsection{9.2. Diseño de Pruebas}\label{sec:cap9_diseno}

% [REDACTAR]

\subsubsection{9.2.1. Identificación y Estructura de Procedimientos}\label{sec:cap9_procedimientos}

Para aplicar el método de caja negra se diseñó una estructura de datos única, que se utiliza para documentar cada caso de prueba. Cada caso de prueba corresponde a un caso de uso del sistema y se registra en una tabla con los siguientes campos, que combinan texto e imágenes de la ejecución:

\begin{itemize}
    \item \textbf{Nombre del caso de prueba:} nombre completo del caso de uso sobre el cual se realiza la prueba, con su identificador.
    \item \textbf{Descripción:} breve resumen del caso de uso evaluado, destacando los aspectos funcionales que se van a verificar.
    \item \textbf{Precondición:} estado o condición que debe cumplirse antes de ejecutar la prueba. Garantiza que el entorno y los datos iniciales sean adecuados para la evaluación, por ejemplo, la sesión de un rol determinado o la existencia de una convocatoria abierta.
    \item \textbf{Datos de entrada:} información que requiere el caso de uso para producir resultados significativos para el usuario, como los valores ingresados en un formulario o los parámetros de una consulta.
    \item \textbf{Condición de ejecución:} restricción que debe cumplirse para demostrar el correcto funcionamiento, cuando el caso de prueba la presenta. Si no existe, se indica que no aplica.
    \item \textbf{Captura de imagen inicial:} captura de pantalla del estado del sistema antes de ejecutar los pasos de la prueba.
    \item \textbf{Pasos:} secuencia numerada de acciones que realiza el probador sobre el sistema, desde el estado inicial hasta obtener el resultado.
    \item \textbf{Resultado esperado:} resultado previsto al finalizar la prueba, ya sea exitoso o fallido. En este punto se establece cómo debe comportarse el sistema según lo especificado.
    \item \textbf{Resultado obtenido:} transcripción de lo que respondió el software al ejecutar la prueba, para su validación o para su corrección si fuese necesario.
    \item \textbf{Captura de imagen final:} captura de pantalla del estado del sistema luego de ejecutar los pasos, como evidencia del resultado obtenido.
    \item \textbf{Resultado de prueba:} veredicto de la prueba, «Correcto» cuando el resultado obtenido coincide con el esperado e «Incorrecto» en caso contrario.
\end{itemize}

La Tabla~\ref{tab:pru_plantilla} presenta la plantilla de esta estructura, que se completa para cada caso de uso en las pruebas de caja negra.

\begin{table}[H]
\caption{Plantilla de caso de prueba de caja negra}
\label{tab:pru_plantilla}
\centering
\footnotesize\linespread{1}\selectfont
\renewcommand{\arraystretch}{1.3}
\begin{tabular}{>{\raggedright\arraybackslash}p{4.5cm}>{\raggedright\arraybackslash}p{10.5cm}}
\toprule
\textbf{Campo} & \textbf{Contenido} \\
\midrule
Nombre del caso de prueba & CU-XX. Nombre completo del caso de uso \\
Descripción & Resumen del caso de uso y de los aspectos funcionales a evaluar \\
Precondición & Estado o condición previa a la ejecución \\
Datos de entrada & Datos que se ingresan al sistema \\
Condición de ejecución & Restricción de la prueba, o «No aplica» \\
Captura de imagen inicial & (Captura de pantalla del estado inicial) \\
Pasos & 1. Acción del probador. 2. Acción del probador. \\
Resultado esperado & Comportamiento previsto del sistema \\
Resultado obtenido & Respuesta real del sistema \\
Captura de imagen final & (Captura de pantalla del resultado) \\
Resultado de prueba & Correcto / Incorrecto \\
\bottomrule
\end{tabular}

\smallskip
\noindent\raggedright\textit{Nota.} Elaboración propia.
\end{table}
